\documentclass[nofootinbib,aip,amsmath,amssymb,reprint,twocolumn,10pt]{revtex4-2}
\usepackage{graphicx}% Include figure files
\usepackage{dcolumn}% Align table columns on decimal point
\usepackage{bm}% bold math
\usepackage[utf8]{inputenc}
\usepackage[T1]{fontenc}
\usepackage{mathptmx}
\usepackage{verbatim}
\usepackage{fdsymbol}
\usepackage{sidecap}
\usepackage{placeins}
\usepackage{xcolor}
\usepackage{amsfonts,amsthm,amscd}
\usepackage{mathtools}
\usepackage{hyperref}
\usepackage{bbold}
% \usepackage{float}
\usepackage{tabularx}
\usepackage[ruled,lined]{algorithm2e}
\usepackage{amssymb} 

\newcommand{\llbracket}{\mathopen{[\![}}
\newcommand{\rrbracket}{\mathclose{]\!]}}

\begin{document}

\title{Kinetics of plasticity regulate burst frequency in a Dendronotus iris Swim CPG Model}

\author{James Scully}
\email{jamesjscully@gmail.com}
\affiliation{Neuroscience Institute, Georgia State University,
100 Piedmont Ave., Atlanta, GA 30303, USA.}
\author{Akira Sakurai}
\email{akira.gsu@gmail.com}
\affiliation{Department of Mathematics \& Statistics, Georgia State University, 
25 Park Pl., Atlanta, GA 30303, USA.}
\author{Paul Katz}
\email{pkatz@umass.edu}
\affiliation{Neuroscience Institute, Georgia State University, 
100 Piedmont Ave., Atlanta, GA 30303, USA.}
\author{Andrey L. Shilnikov}
\email{ashilnikov@gsu.edu}
\affiliation{Neuroscience Institute and Department of Mathematics \& Statistics, Georgia State University, \\
100 Piedmont Ave., Atlanta, GA 30303, USA.}

\keywords{locomotion, modeling, half-center oscillator, hco, excitatory, inhibitory, rhythmic, oscillations, bursting, neurons, subnetworks, synaptic, metabotropic, network, hysteresis, neuroscience, fast-slow, adaptability, robustness, resilience}

\date{\today}
\begin{abstract}


The swim central pattern generator (CPG) of Dendronotus iris produces robust rhythmic activity despite variability in its components and external perturbations.
This study examines the contributions of slow synaptic dynamics and the coexistence of multiple oscillatory pathways to the resilience of this network.
We constructed a four-cell model based on experimental data, incorporating slow synaptic processes and fast-slow decomposition.
Preliminary findings indicate that robustness emerges from the combined influence of multiple oscillatory mechanisms rather than reliance on a single dominant pathway.
This model provides a framework for understanding adaptive motor rhythms and their application to biological and synthetic systems.
\end{abstract}

\maketitle
\section{Introduction}
Animal locomotion emerges from neural circuits that must remain effective despite intrinsic variability, changing body size, and fluctuating environmental conditions.
Early work on locusts demonstrated ``fictive locomotion,'' revealing that rhythmic motor output can persist without sensory feedback \cite{Wilson1961}.
This led to the identification of autonomous, rhythmogenic circuits known as central pattern generators (CPGs) that drive locomotion in both invertebrates and vertebrates \cite{MarderCalabrese1996, Kiehn2006}.


Because locomotion underpins critical behaviors like foraging and predator avoidance, CPGs evolve under selective pressure to adapt to environmental variables like temperature and maintain appropriate scaling among body size, locomotion speed, and rhythmic frequency \cite{Marder2023}.
These circuits integrate sensory feedback and descending control signals, relying on dynamical mechanisms across multiple timescales: spikes at milliseconds, bursts shaped over seconds, and homeostatic feedback stabilizing circuits over minutes to days \cite{MarderCalabrese1996}.

Prior research has shown that ion channel degeneracy allows broad variability in conductances among genetically identical individuals to yield similar functional states \cite{Marder2011, OLeary2014}, with homeostatic feedback linked to intracellular calcium stabilizing neuronal activity \cite{Franci2020}.
However, this does not explain how synapses and cells coordinate to produce reliable bursting across a range of frequencies.
We propose that controllability, the ability to produce diverse behaviors by modulating a control variable, is another key factor.
In this study, we demonstrate how the kinetic properties of synaptic enhancement allow burst frequency control in the \textit{Dendronotus iris} swim CPG \cite{SakuraiKatz2016}.

Initially driven by evolutionary questions about locomotor diversity within its clade---including dorsal-ventral swimmers, left-right swimmers, and non-swimmers \cite{Newcomb2012}---our research on the \textit{Dendronotus iris} swim CPG uncovered a novel role for neuroplasticity.
Our modeling indicates that the temporal properties of synaptic enhancement, rather than its magnitude or static strength, confer precise control over the network's bursting frequency and pattern [neuroplasticity citation].
Specifically, the kinetics of synaptic enhancement provide network-level adaptability in this system \cite{SakuraiKatz2016}.

Looking ahead, this work highlights the critical role of synaptic plasticity kinetics in governing neural circuit dynamics [neuroplasticity citation].
These properties may serve as a general mechanism for controllability across diverse systems, though the biological processes regulating them remain underexplored.
Synaptic enhancement spans facilitation, augmentation, and post-tetanic potentiation, with augmentation matching the timescale of bursting in the \textit{Dendronotus iris} CPG [neuroplasticity citation].
This could involve co-regulation of enhancement forms or selective modulation of augmentation, possibly via synaptotagmin isoforms [neuroplasticity citation].

The paper is organized as follows: First, we describe the swim CPG network assembly, detailing neuron and synaptic interaction modeling.
Second, we validate the network by comparing simulated voltage traces with empirical data from laboratory experiments involving current injections and perturbations.
Finally, we explore how burst frequency is precisely controlled through synaptic enhancement kinetics, demonstrating the pivotal role of intermediate timescale dynamics in adaptive locomotion.


\section{Results: Body size swim patterns, fact checks or verification of the CPG models}


\begin{figure*}[tbp]
\centering
\includegraphics[width=.75\textwidth]{../"corel draw"/collage_a.jpg}
\caption{
%Three vertical panels in the collage to illustrate a wide range of swim patterns generated by the swim CPG circuit to match the maturing animal size. (A) Swim CPG circuitry with fore core interneurons Si2L/R and Si3L/R  driven by ``sensory'' interneurons Si1s.  (B) Voltage traces (top)   along with the rates [probabilities[ of neurotransmitter released by the swim CPG core (four) interneuron (SiN) models.  (C) The weight (size) of the  increases nearly 10 times over the animal life span: from 30 grams up to 300 grams, while its swim period slows down from 2-3 seconds in juveniles up to 12-14 seconds in aging species, see the voltage traces in panel B.
%
A three-panel collage illustrating diverse swim patterns generated by the swim CPG circuit, adapting to the maturing size of Dendronotus iris. (A) Swim CPG circuitry featuring core interneurons Si2L/R and Si3L/R, driven by the sensory interneurons Si1s.
(B) Voltage traces (top) and neurotransmitter release probabilities for the four core swim CPG interneurons (SiN). 
(C) As {\em Dendronotus Iris} grows from 30g to 300g over its lifespan, as shown in the voltage traces in panel~B. its swim period slows down from 2--3s in juveniles to up to 12--14 seconds in older adults, see the voltage traces in panel~B.}
\label{fig:group_fig5}
\end{figure*}

\begin{figure*}[tbp]
\centering
\includegraphics[width=.55\textwidth]{../"corel draw"/collage3a.jpg}
\caption{
Anti-phase left–right episodes within a swim pattern, produced by the Dendronotus iris CPG under control conditions, arise from twisted pairwise modules (Si3R $\leftrightarrows$ Si2L and Si3L $\leftrightarrows$ Si2R) that are unidirectionally coupled via excitatory ($\blacktriangledown$) and inhibitory ($\bullet$) synapses. These modules alternate in initiating the onset of network bursting, an effect further reinforced by two reciprocally inhibitory HSOs, consisting respectively of Si3R–Si3L and Si2R–Si2L pairs. Compare with the original Fig.1~from the experimental study~\cite{sakurai2022bursting}.
}
\label{fig:group_fig5}
\end{figure*}





\subsection{Resilience to pulse perturbations. Fact check list I}

\begin{figure*}[tbp]
\centering
\includegraphics[width=.49\textwidth]{../jack_figs/5a_als.jpg}~~~~~
\includegraphics[width=.49\textwidth]{../jack_figs/5b_als.jpg}
\includegraphics[width=.49\textwidth]{../jack_figs/5c_als.jpg}~~~~~
\includegraphics[width=.49\textwidth]{../jack_figs/5e_als1.jpg}
\includegraphics[width=.49\textwidth]{../jack_figs/5f_als1.jpg}~~~~~
\includegraphics[width=.49\textwidth]{../jack_figs/5g_als.jpg}
\caption{Comparison of specific, distinct responses of the swim CPG to targeted perturbations, produced by injecting hyperpolarizing and depolarizing current pulses (red vertical arrows) into selected interneurons. This comparison validates the CPG model against its biological counterpart and demonstrates the resilience of both systems in restoring and maintaining normal swim patterns.
Fact-check list~I: Compare these six modeling panels with their corresponding experimental originals in Figs. 5A–C and E–G of Ref. \cite{SakuraiKatz2016}. For empirical evidence of resilience, see Fig.~5 in Ref. \cite{SakuraiKatz2017}. For the PIR proof, see Fig. 5 in the other experimental Ref.~\cite{sakurai2022bursting}.}.
\label{fig:group_fig5}
\end{figure*}

\subsection{Resilience to pulse perturbations. Fact check II}


\begin{figure*}[tbp]
\centering
\includegraphics[width=.37\textwidth]{../jack_figs/7a_als1.jpg}~~~~~
\includegraphics[width=.37\textwidth]{../jack_figs/7b_als2.jpg}\\
\includegraphics[width=.37\textwidth]{../jack_figs/7c_als.jpg}~~~~~
\includegraphics[width=.37\textwidth]{../jack_figs/7d_als1.jpg}
\caption{Similar responses of the swim CPG during control swims to various temporal perturbations—produced by applying hyperpolarizing or depolarizing current pulses (red vertical arrows) to targeted interneurons—allow us to validate the CPG model against its biological counterparts.
Fact-check list~II: compare the four modeling panels with their corresponding experimental originals shown in Figs. 7Ai–-Di of Ref.~\cite{SakuraiKatz2016}.
}
\label{fig:group_fig7}
\end{figure*}


\subsection{Si1 controls swim initiation and termination. Fact check III}



\begin{figure*}[tbp]
\centering
\includegraphics[width=.55\textwidth]{../jack_figs/s1_Fig2A_drives_cpg.jpg}\\
\includegraphics[width=.55\textwidth]{../jack_figs/s1_Fig3A_drives_cpg1.jpg}\\
%\includegraphics[width=.75\textwidth]{../jack_figs/s1_Fig3B_drives_cpg_als}
\includegraphics[width=.55\textwidth]{../jack_figs/s1_drives_CPG_Fig3B_No_NM.jpg}
\caption{Interneuron Si1 excites Si3 neurons, thereby initiating and terminating the entire rhythmically bursting CPG (in the absence of neuromodulation).
Fact-check list~II to compare:
i) Panel A with its experimental original — Fig. 2A from Ref. \cite{sakurai2019command};
ii) Panels B and C with their corresponding originals — Figs. 3A and 3B from Ref. \cite{sakurai2019command}, respectively.}
\label{fig:group_fig5}
\end{figure*}


\subsection{Role of neuromodulation and its modeling. Burst frequency control}


\section{Methods: cellular and synaptic models}



\begin{figure}[tbp]
\centering
\includegraphics[width=.25\textwidth]{../jack_figs/ZOO_synapses.jpg}\\
\includegraphics[width=.22\textwidth]{../jack_figs/S1_synapses1.jpg}
\includegraphics[width=.22\textwidth]{../jack_figs/S1_synapses2.jpg}
\caption{Fact check IV: Compare the panels B and C with the original neurophysiological data — Fig. 7 in Ref. \cite{sakurai2022bursting}.
Three positive current pulses of increasing amplitude were applied to Si2R to calibrate the $\alpha$-synapse and logistic synapse models (bottom trace), as well as the IPSP responses in the contralateral Si2L (top trace). Two positive current pulses injected into Si2R triggered accumulating IPSP responses in Si2L (top trace) and elicited fast and slow echoes in Si3L. These echoes arise from the fast gap junction and delayed slow inhibition mediated by the logistic synapse between Si2R and Si3L, respectively (see the swim CPG circuitry in Fig.~1).}
\label{fig:group_fig5}
\end{figure}


\section{CPg assembly line}


\begin{figure*}[tbp]
\centering
\includegraphics[width=.4\textwidth]{../jack_figs/assembly_step1.jpg}
\includegraphics[width=.4\textwidth]{../jack_figs/assembly_step2.jpg}\\
\includegraphics[width=.4\textwidth]{../jack_figs/assembly_step3.jpg}
\includegraphics[width=.4\textwidth]{../jack_figs/assembly_step4.jpg}\\
\includegraphics[width=.4\textwidth]{../jack_figs/assembly_step5.jpg}\\
\caption{CPG assembly line in five sequential steps (order is arbitrary):
(A) Four uncoupled interneurons, with Si3 neurons firing at a higher rate due to excitatory drive from Si1 neurons.
(B) Si3 neurons provide fast excitatory drive to ipsilateral Si2 neurons.
(C) Si2 neurons deliver inhibitory feedback onto Si3 neurons.
(D) Reciprocally inhibitory synapses are activated between contralateral Si3 neurons.
(E) Reciprocally inhibitory synapses are activated between contralateral Si2 neurons, along with gap junctions between Si3 neurons, which together generate stable fast network-level bursting in the swim CPG model of {\em Dendronotus iris}.}
\label{fig:group_fig5}
\end{figure*}


The central objective of this study is to develop and validate a computational model of the \textit{Dendronotus iris} swim central pattern generator (CPG) that enhances network robustness through a four-cell circuit.
This was achieved by constructing the model stepwise, beginning with well-characterized two-cell subnetworks and integrating them into a larger network.
A systematic approach ensured biological plausibility, dynamic stability, and alignment with empirical data.
Detailed analysis of the neuronal, synaptic, and subnetwork models, including fast-slow decomposition and two-cell oscillatory subnetwork construction, is provided in \cite{scully2024pairing}. We summarize the essential details here for convenience.

\subsection*{Model of the swim interneuron (SiN) }
We adapted a conductance-based swim interneuron (SiN) model from Plant's 1981 framework \cite{plant81, scully2024pairing}, modifying it with an additional fast $I_h$ current to represent the cellular components of the \textit{Dendronotus iris} swim CPG.
The governing equation is:
\begin{equation}
C_m \frac{dV}{dt} = -I_{Na} - I_{Ca} - I_K - I_\text{leak} - I_h - I_{KCa},
\end{equation}
where $C_m$ (1 $\mu$F/cm$^2$) is the membrane capacitance, $V$ is the membrane voltage, and ionic currents include:
%
\begin{itemize}
    \item \textbf{Fast inward current}: A combined sodium-calcium current with instantaneous activation and slow inactivation, akin to Hodgkin-Huxley, driving action potentials.

    \item \textbf{Slow inward current}: A TTX-resistant current with slow activation, contributing to bursting periodicity.
    \item \textbf{Fast outward current}: A potassium current, Hodgkin-Huxley-like, for re-polarization.
    \item \textbf{Slow outward current}: A calcium-sensitive potassium current, tied to slow intracellular calcium buildup, shaping burst duration.
    \item \textbf{Leak current}: A passive inward current for baseline dynamics.
    \item \textbf{Fast hyperpolarization-activated current}: $I_h$, added to enhance dynamics at hyperpolarized voltages.
\end{itemize}

The original Plant model includes the following fast currents: the inward sodium and calcium ($I_{I}$), the outward potassium ($I_K$), necessary for the spike generation, a depolarizing h-current ($I_h$), along with the generic ohmic leak ($I_{leak}$) current. The gradual spike frequency adaptation and post-inhibitory rebound in the model are due to the slow dynamics of two currents: the TTX-resistant inward sodium and calcium current ($I_{\rm T}$) and outward calcium-sensitive potassium current ($I_{KCa}$). 
%
\begin{gather}
C_{m} {V}^\prime= - I_{I} - I_{K} - I_{leak} - I_{h} - I_{T} - I_{KCa}, \label{eq:voltage} ~\\ 	
h^\prime = \frac{h_{\infty}(V)-h}{ {\tau_{h}(V)} }, \qquad 
n^\prime = \frac{n_{\infty}(V)-n}{\tau_{n}(V)}, \label{eq:hngating} \\
y^\prime= \frac{1}{2} \left [ \frac{1}{1+e^{10(V-50)}}-y \right ]/ \left [ 7.1+\frac{10.4}{1+e^{(V+68)/2.2}} \right ] \label{eq:hgating}
\end{gather}
with the membrane capacitance $C_m=1$. Here the dynamic variables are the membrane voltage $V(t)$, the gating probabilities $h(t)$, $n(t)$, and $y(t)$. This ODE system describing the R15 Plant burster includes a fast subsystem to generate repetitive tonic-spiking or quiescent activity, depending on the level of drive from its slow subsystem. This spike-generating machinery of the fast subsystem consists of the four aforementioned currents, which are governed by the following equations:
% 
\begin{align}
I_{I} &= g_{I}\,h\, m^{3}_{\infty}(V)\,(V-E_{I}), \label{eq:icurrent}\\ 
I_{K} &= g_{K} \, n^{4}\, (V-E_{K}), \label{eq:kcurrent} \\ 
I_{h} &= g_{h}\,\frac{y\,(V-E_{h}) }{\left (1+e^{-(V-63)/7.8} \right )^3 } \label{eq:hcurrent}, \\
I_{leak} &= g_{L} (V-E_{L})
\end{align}
Here, the activation of the inward sodium current is assumed instantaneous and therefore described by an analytical (sigmoidal) equation $m_{\infty}^3(V)$, rather than a corresponding fast ODE. The third current describes a depolarizing $h$-current that activates as the voltage drops down below $-50$mV, see Eq.~(\ref{eq:hgating}) above. \footnote{The Plant model generates bursts with deeply hyperpolarized quiescent sags, see Fig.~\ref{fig:Plant_burster}(B). To flatten/rectify those sags, the $h$-current was additionally introduced in the interneuron model, and we would like to thank to D.~Terman for the suggestion. This current also prevents unrealistic levels of hyperpolarization in response to negative external currents applied to the neuron model. The use of $h$-current, though practical, does not appear to play a critical role in the swim CPG models.} 

Recurrent alternation between fast spike trains and slow quiescent episodes in the endogenous Plant burster is reciprocally modulated by the slow inward TTX-resistant $\rm Na^{+}$-$\rm Ca^{2+}$ current and slow outward $\rm Ca^{2+}$ activated $K^+$ given, respectively, by 
%
\begin{align} 
I_{T} &= g_{T} x (V-E_{I}), \label{eq:ttxcurrent} \\
I_{KCa} &= g_{KCa}\frac{[\rm Ca]}{0.5+[\rm Ca]}(V-E_{K}) \label{eq:cacurrent} 
\end{align}
with two dynamical variables: the calcium concentration $[\rm Ca(t)]$ and a voltage gated probability $x(t)$, which are governed by following coupled slow system:  
%
\begin{align}
x^\prime &= \frac{1}{\tau_x}\left [ \frac{1}{1+e^{-0.15(V+50- \Delta_{x})}}-x \right ], &\tau_x &\gg 1, \label{SlowSub1} ~\\
[\rm Ca]^\prime  &=  \rho \left ( K_{c} x (E_{Ca}-V + \Delta_{\rm Ca})-[\rm Ca] \right  ), &\rho &\ll 1,  \label{SlowSub2}  
\end{align}
where $\Delta_x$ and $\Delta_{\rm Ca}$ are bifurcation parameters introduced to control slow dynamics of the SiN-model; we will discuss their action below. The reader can find the detailed description of biophysical parameters in the Appendix ~\hyperref[sec:appendix1]{I} below.

The time scales of the state variables: $V(t)$, probabilities $h(t)$, $n(t)$, $y(t)$ and $x(t)$ and calcium concentration $[\rm Ca](t)$ that gate the indicated cellular currents in the model are presented in Fig.~\ref{fig:currents}. Note that depending on the chosen value of the time constant $\tau_x$, the average rate of change of the gating $x$-variable can be altered: it  can be visually as fast (at $\tau_x=100$ sec$^{-1}$) as that of the fast subsystem, or can be slowed down to match the rate of change of the $[\rm Ca]$-variable at larger values such as $\tau_x=273$ sec$^{-1}$ also used in this study. Nonetheless, the $x$-dynamics is included in the slow subsystem for a few reasons. First, for the sake of historical continuity, $x(t)$ was treated as a slow variable in the original slow-fast dissection analysis of this model proposed in Ref.~\cite{RL87}. Second, the $x$-dynamics does not contribute (but regulate) to the spike generation, which is solely due to the fast sodium-calcium and potassium currents. Instead its reciprocal nonlinear interaction with the [\rm Ca]-dynamics is the key factor that provides the Plant model with a slow hysteresis necessary for the onset of endogenous bursting composed of alternating fast spike trains and quiescent episodes. 

The parameter $\Delta { [\rm Ca]}$ controls the intrinsic excitability of the neuron. Decreasing it is equivalent to increasing its excitability. Specifically, the $\Delta_x$-parameter represents a deviation from the voltage value $-50$mV at which the TTX-resistant $\rm Na^{+}$-$\rm Ca^{2+}$ current is half-activated, see   Eq.~(\ref{SlowSub1}) and the corresponding activating function $x_{\infty}(V)=1/ \left (1+e^{-0.15(V+50-\Delta_{x})} \right )=1/2$. 
The parameter $\Delta _{x}$ changes how the cell responds qualitatively to synaptic drive. When $\Delta _{x}$ is 0, the model has an endogenous bursting regime. Decreasing $\Delta _{x}$ to -4mV eliminates this endogenous bursting. Specifically, the $\Delta_{\rm Ca}$-parameter  introduced in Eq.~(\ref{SlowSub2}) shifts the calcium reversal potential from a hypothetically high value $+140$mV. This voltage is too high to be measured experimentally, so changing the reversal potential is biologically plausible.





Slow variables, intracellular calcium concentration $[\rm Ca]$ and gating variable $x$, evolve as:
\begin{equation}
\frac{dx}{dt} = \frac{x_\infty(V) - x}{\tau_x},\quad 
\frac{d[\rm Ca]}{dt} = \rho \left( K_c x (E_{Ca} - V + \Delta { [\rm Ca]}) - [\rm Ca] \right),
\end{equation}
where $x_\infty(V)$ and $\tau_x$ (ms) are the voltage-dependent activation and time constants, and $\rho$ (ms$^{-1}$) and $K_c$ (mM$^{-1}$) control calcium influx and decay.




The full description of the model equations are given in Ref.~\cite{scully2024pairing}, see also the Appendix~I below.


\subsection*{Synaptic Models}

Synaptic dynamics are modeled using a phenomenological approach detailed in \cite{scully2024pairing}, which balances realism and tractability.
Two synapse types were implemented:
\begin{enumerate}
    \item \textbf{Fast synapses}: Modeled with a gating variable $S$ that rises quickly with presynaptic voltage and decays at a fixed rate ($\alpha = 10$ ms$^{-1}$, $\beta = 5$ ms$^{-1}$), following the $\alpha$-synapse formulation \cite{destexhe1994synthesis}.
    \item \textbf{Slow synapses}: Include a nonlinear term for synaptic accumulation (minima above zero) and high-pass filtering of low-frequency inputs.
\end{enumerate}
Full equations are in \cite{scully2024pairing} or the appendix.

The synapses in the \textit{Dendronotus} swim CPG may be either fast or slow. Slow synapses are of particular interest, as they occur on the same timescale as network oscillations and have been reported in the swim CPGs of the sea slugs \cite{NW02, WNT02, TW05, SNLK11}. In order to be able to accurately control the frequency response of the synaptic gating variable we introduce a new {\em logistic} model that can describe both delayed and slow synaptic summation, which can act as an efficient high-pass filter. We briefly compare the logistic synapse's properties with similar synapses from literature, namely the $\alpha$-synapse, \cite{Wang-Rinzel-92, Rall67}, higher order kinetics \cite{destexhe1998kinetic,DESTEXHE-CONTRERAS-SEJNOWSKI-STERIADE-94, Golomb96}, and a dynamic synapse~\cite{Hill2001}.


The relationship between pre-synaptic frequency and neurotransmitter release is computed for each synapse from trajectories of a pre-synaptic cell coupled to a post-synaptic cell. The frequency of the pre-synaptic cell is increased over time by slowly decreasing the $\Delta_{\rm{Ca}}$ parameter. For each spike over this trajectory, the frequency is calculated as the reciprocal of the inter-spike interval, and the average neurotransmitter rate, $\angle S \rangle$ is computed from the time-average of the synaptic gating variable $S$ over the same spike. The time series demonstrating this comparison is pictured in Fig ~\ref{fig:five_panels}, and the curves showing the relationship between pre-synaptic frequency and $<S>$ are shown in Fig ~\ref{fig:freq}.

The synaptic current $I_{syn}$ in the post-synaptic neuron is modeled as follows:  
\begin{equation}
\quad I_{syn} = g_{syn}\, S(t) \, (V_{post}(t)-E_{rev}), \label{Eqsyn}
\end{equation}
where $g_{syn}$ is a maximal conductance, $S(t) $ is a synaptic gating variable, $V_{post}$ is the membrane voltage in the post-synaptic neurons, and $E_{rev}$ is a synaptic reversal potential, which can be set at +40mV or -80mV for excitatory and inhibitory synapses, respectively. The sigmoid  function $f_{\infty}(V)$ is defined as
\begin{equation}
f_{\infty}(V) = \frac{1}{1+e^{-k(V_{pre}-\Theta_{syn})}},  \label{finf}
\end{equation}
where the constant $k$ determines its derivative at $f_{\infty}(V) = 0.5$, the inflection point where $V=\Theta_{syn}$. Here $\Theta_{syn}$ is treated as the synaptic threshold typically set around $+20$mV in the middle of spikes, between the spike threshold around $-40$mV (sodium channels opens) and the spike peak $+40$mV. Values for $k$ are typically set somewhere in a range $\left[0.5,\, 50\right]$, and as the value of $k$ is set large,  the function becomes a continuous approximation of the Heaviside step function, where the function $f_\infty$ remains close to zero as long as $V_{pre}(t) < \Theta_{syn}$, and quickly jumps to 1 after the voltage in the presynaptic neurons exceed the synaptic threshold.

\paragraph{$\alpha$-synapse.} Adopting the phenomenological approach taken by Hodgkin and Huxley, Wang and Rinzel~\cite{Wang-Rinzel-92} proposed what is now commonly referred to as an $\alpha$-synapse. Its activation is meant to mimic the profile of an $\alpha$-function given by $t^p\, e^{-t}$, where $p$ is a positive integer. The idea of such $\alpha$-synapses is rooted in the pioneering computational work by W.~Rall~\cite{Rall67}, who studied and modeled various aspects of synaptic potentials.
 
The dynamic equation describing the rate of change of the $S(t)$-variable in the $\alpha$-synapse is given by
\begin{equation}
S'(t) = \alpha (1-S) f_{\infty}(V_{pre}) - \beta S \label{Salpha}
\end{equation}
with some positive  $\alpha$ and $\beta$ constants. Here, $f_{\infty}$ function with the presynaptic membrane potential $V_{pre}$: 
\begin{equation}
S(t) = f_{\infty}(V_{pre}(t)). \label{ftm}
\end{equation}
 Such fast synapses happened to be useful for understanding synchronized neuronal activity, including bursting with inhibitory synapses, and multistability in small, weekly coupled neural networks~\cite{prl08, Shilnikov2008b, Jalil-Belykh-Shilnikov-10, Wojcik2011a,  Wojcik2014, schwabedal2016qualitative,kelley20202,collens2020dynamics, pusuluri2020computational}. Moreover, this assumption has the benefit of simplifying the network dynamics, facilitating both analysis and simulation \cite{jalil2013,Schwabedal-Neiman-Shilnikov-14,lodi2019design}. 


So, when $V(t)$ is below the synaptic threshold $\Theta_{syn}$ and hence $f_{\infty}=0$, then the synaptic probability $S(t)$ exponentially decreases to zero as $e^{-\beta t}$. During an action potential as long as $V(t) \ge \Theta_{syn}$, then $S(t)$  is approaching the equilibrium state ($\alpha/(\alpha+\beta$)) exponentially fast as $e^{-(\alpha+\beta)t}$. Note that Eq.~(\ref{Salpha}) is sometimes referred to as an $\alpha-\beta$-synapse with the first-order kinetics \cite{destexhe1994synthesis}.

If $\alpha$ and $\beta$ values are on the same order of magnitude, say, 0.1, then the $\alpha$-synapse is as fast as an FTM \cite{Alacam2015,Jalil-Belykh-Shilnikov-12,baruzzi2021towards, baruzzi2020generalized}. Decreasing $\beta$ by one order of magnitude $\sim 0.01$ makes the synapse sufficiently slow and stronger as it can accumulate and demonstrate the pronounced summation with monotonically increasing $S(t)$ on average, see Fig.~\ref{fig:five_panels}A3 and B3 (grey lines), in response to trains of fast spikes in the pre-synaptic neuron. Typical values for the $\alpha$-synapses in the swim CPGs are $\alpha = 0.01$, while $\beta$ is set in the range $\left[0.001,\, 0.0005\right]$ for slow synapses, and $\left[0.01,\, 0.1\right]$ for fast ones. Decreasing proportionally $\alpha$ and $\beta$ makes the time progression of $S(t)$ smoother with smaller amplitude variations.      

Most slow synapses act through complex sequences of nonlinear interactions including secondary messenger cascades, which can delay the neurotransmitter release in pre-synaptic neurons or their binding with neurotransmitter receptors in post-synaptic ones. In such cases it is sometimes desirable to introduce a dynamic delay or some low spike frequency filter without explicitly modeling every stage of postsynaptic intracellular signal transduction in terms of state variables. For this purpose, we introduce a new synapse model, which we call the logistic synapse, named after its similarity to the logistic model of population growth, where the probability $S(t)$ is governed by a single ODE:  
\begin{equation} \label{sigmoid}
S'(t) = \alpha\, S\,(1-S)\, f_{\infty}(V_{pre}) - \beta (S-S_0), \\
\end{equation} 
where a constant $0 \le S_0 \le 10^{-3}$ can be viewed as some spontaneous neurotransmitter release rate on average in the pre-synaptic neuron. The key feature of this modeling approach is the term $S\,(1-S)$  which provides: (i) means to control a latency of the synapse in the initial ascending phase of summation, and hence (ii) a logistic or sigmoid-like shape of the $S(t)$-variable determining the proportion of open channels. With this logistic model, we obtain a desired nonlinear dependence of the strength of the synaptic current on the spike frequency in the pre-synaptic neurons to match the experimental studies on the swim CPGs. In addition, with the logistic synapse acting as a high-pass filter we can further explore the role of such time-varying synaptic coupling on the given CPG rhythmogenesis and its robustness.




\subsection*{Four-Cell Network Assembly}
We constructed the four-cell network using validated two-cell or pairwise subnetworks 
such as half-center oscillators 
(HCOs) and  excitatory-inhibitory (EI) modules introduced in Ref.~\cite{scully2024pairing}.
 
 
Synaptic kinetics for mutual inhibition matched HCO results, and E/I interactions were similarly aligned.

We preserved cellular parameters and synaptic kinetics, adjusting only the endogenous excitability of Si3 (swim interneuron 3) by shifting its voltage threshold $\Delta [\rm Ca]$ by 5mV to enhance spiking.
Calibration relied on symmetrically tuning synaptic conductances, following these steps:\\
%
\begin{enumerate}
    \item \textbf{Configuration.} Substituted pre- and postsynaptic voltage variables to reflect the known \textit{Dendronotus iris} circuit \cite{SakuraiKatz2016}.
    \item \textbf{Initialization.} Set all cells (Si1--Si4) to a quiescent state and synaptic conductances to zero.
    \item \textbf{Drive.} Applied synaptic drive to Si3 until the excitatory postsynaptic $S$
    accumulated with minima above 0.1, saturating over 0.5--3 s (1/4 to 1/2 the burst period
     of 2--12s, varying with animal size).
    \item \textbf{Excitation.} Increased excitatory conductance from Si3 to contralateral Si2 until inhibitory $S$ (Si2-to-Si2) accumulated with minima above 0.1 over 0.5--3 s.
    \item \textbf{Slow Inhibition.} Raised E/I pair inhibitory conductance until Si3 frequency dropped by $50\%$ over 1 s.
    \item \textbf{Fast Inhibition.} Elevated inhibitory conductances for both inhibitory HCOs (Si1--Si2 and Si3--Si5) until a stable burst rhythm emerged, visually resembling empirical voltage traces \cite{SakuraiKatz2016}.
    \item \textbf{Fine-Tuning.} Adjusted excitability slightly (e.g., $\Delta [Ca] < 2$ mV) to set a 30\% duty cycle, verified as extrinsic by resetting conductances to zero.
\end{enumerate}

\textcolor{red}{TODO: Are we using Si1 or Iapp for Si3 drive?}

\subsection{Validation from voltage perturbation}


\label{fig:supp} 
\section*{Appendix I: a basic Hodgkin-Huxley based model of the swim CPG interneuron (SiN)} \label{sec:appendix1}
The dynamics of the membrane potential, $V$, is governed by the following equation:
\begin{equation}
 	C_{m} {V}^\prime = -I_{I} - I_K - I_{T}  - I_{KCa} -  I_{h}  - I_{leak}  - I_{syn}.  \label{a1}
 \end{equation}
 The fast inward sodium and calcium $I_{I}$-current is given by   
\begin{equation}
I_{I}=g_{I}\,h\, m^{3}_{\infty}(V)(V-E_{I}), 
 \end{equation}
with the reversal potential $E_{I}=30$mV and the maximal conductance value $g_{I}=4$nS, and  
\begin{align}
 m_{\infty}(V) &= \frac{\alpha_{m}(V)}{\alpha_{m}(V) +\beta_{m}(V)}, \qquad  \\
 \alpha_{m}(V) &= 0.1 \frac{50-V_s}{-1+e^{(50-V_s)/10}}\ , \\ 
 \quad \beta_{m}(V) &= 4 e^{(25-V_s)/18}  
\end{align}
 the dynamics of its inactivation  gating variable $h$ is given by      
\begin{align}
 {h}^\prime &=  \frac{h_{\infty}(V)-h}{ {\tau_{h}(V)} }, \quad \mbox{where} \quad \\
h_{\infty}(V) &= \frac{\alpha_{h}(V)}{\alpha_{h}(V) +\beta_{h}(V)}\quad \mbox{and} \\
 \quad \tau_{h}(V) &= \frac{12.5}{\alpha_{h}(V) +\beta_{h}(V)} 
\end{align}
 \begin{align}
 \alpha_{h}(V) &= 0.07 e^{(25-V_s)/20}, \\
 \beta_{h}(V) &= \frac{1}{1+e^{(55-V_s)/10}}, \quad \mbox{and} \quad \\
 V_s &= \frac{127V+8265}{105}{\rm mV}.
 \end{align} 
  The fast potassium $I_K$-current is given by the equation  
\begin{equation}
I_K=g_{K} n^{4}(V-E_{K}),  
\end{equation}
with the reversal potential $E_{K}=-75$mV and the maximal conductance set as $g_K=0.3$nS. The dynamics of inactivation gating variable are described by 
\begin{align}
{n}^\prime &= \frac{n_{\infty}(V)-n}{\tau_{n}(V)}, \quad \mbox{with} \quad \\
n_{\infty}(V) &= \frac{\alpha_{n}(V)}{\alpha_{n}(V) +\beta_{n}(V)} \quad \mbox{and} \quad \\
\tau_{n}(V) &= \frac{12.5}{\alpha_{n}(V) +\beta_{n}(V)}, \quad \mbox{where} \quad \\
\alpha_{n}(V) &= 0.01 \frac{55-V_s}{e^{(55-V_s)/10}-1} \quad \mbox{} \quad \\
\beta_{n}(V) &= 0.125 e^{(45-V_s)/80}.
\end{align}	
The rapidly depolarizing h-current is given by 
\begin{equation}
I_{h}=g_{h} \frac{y (V-E_{h})}{(1+e^{-(V-63)/7.8)})^3}, 
\end{equation}
with $E_{h}=70$mV, and $g_h=0.0006$nS; the dynamics of its $y$-activation probability is described by 
 \begin{equation}
{y}^\prime = \frac{1}{2} \left [ \frac{1}{1+e^{10(V-50)}}-y \right ]/ \left [ 7.1+\frac{10.4}{1+e^{(V+68)/2.2}} \right ], 
\end{equation}	
whereas the leak current is given by 
\begin{align}
I_{leak} &=g_{L} (V-E_{L}), 
\end{align}
with $E_{L} =-40{\rm mV}$, and $g_{L}=0.003{\rm nS}$. The sub-group of slow currents in the model includes the TTX-resistant sodium and calcium $I_{T}$-current given by   \begin{align}
 I_{T} &= g_{T} x (V-E_{I}),
\end{align}
with $E_{I} =30{\rm mV}$ and $g_{T}=0.01{\rm nS}$, where the dynamics of its slow activation variable are described by  
 \begin{align}
{x}^\prime &= \frac{x_{\infty}(V)-x}{\tau_{x}}, \quad \mbox{where} \\ 
\quad x_{\infty} &= \frac{1}{1+e^{-0.15(V+50-\Delta_x)}},
\end{align}\label{eqttx} 
with the time constant $\tau_{x}$ set as 100 or 235ms in this study. The slowest outward $Ca^{2+}$ activated $K^+$ current given by  
\begin{equation}
I_{KCa} =g_{KCa}\frac{[Ca]_i}{0.5+[Ca]_i}(V-E_{K})
 \label{IKCa}
\end{equation}
with  $E_{K}=-75\mbox{mV}$, and $g_{KCa}=0.03$, where the dynamics of intracellular calcium concentration is governed by 
\begin{equation}
{[Ca]}^\prime  =  \rho \left ( K_{c}\, x\, (E_{Ca}-V + \Delta_{Ca})-[Ca] \right  )  \label{eqCa}
\end{equation}
with the Nernst reversal potential  $E_{Ca}=140$mV, and small constants $\rho= 0.0003{\rm ms}^{-1}$ and  $K_c=0.0085{\rm mV}^{-1}$.

\section*{References}
\begin{thebibliography}{10}

  
\bibitem{Wilson1961}
DM~Wilson.
\newblock The central nervous control of flight in a locust.
\newblock {\em Journal of Experimental Biology}, 38(2):471--490, 1961.

\bibitem{MarderCalabrese1996}
E~Marder and RL~Calabrese.
\newblock Principles of rhythmic motor pattern generation.
\newblock {\em {Physiological Reviews}}, {76}({3}):687--717, {1996}.

\bibitem{Kiehn2006}
O~Kiehn.
\newblock Locomotor circuits in the mammalian spinal cord.
\newblock {\em Annual Review of Neuroscience}, 29:279--306, 2006.

\bibitem{Marder2023}
E~Marder.
\newblock Individual variability, statistics, and the resilience of nervous systems of crabs and humans to temperature and other perturbations.
\newblock {\em eNeuro}, 10(11), 2023.

\bibitem{Marder2011}
E~Marder.
\newblock Variability, compensation, and modulation in neurons and circuits.
\newblock {\em Proceedings of the National Academy of Sciences USA}, 108(Suppl 3):15542--15548, 2011.

\bibitem{OLeary2014}
T~O'Leary, AH~Williams, A~Franci, and E~Marder.
\newblock Cell types, network homeostasis, and pathological compensation from a biologically plausible ion channel expression model.
\newblock {\em Neuron}, 82(4):809--821, 2014.

\bibitem{Franci2020}
A~Franci, T~O'Leary  and J~Golowasch.
\newblock Positive dynamical networks in neuronal regulation: How tunable variability coexists with robustness.
\newblock {\em IEEE Control Systems Letters}, 4(4):946--951, 2020.


\bibitem{SNLK11}
A~Sakurai, JM~Newcomb, JL~Lillvis, and PS~Katz.
\newblock Different roles for homologous interneurons in species exhibiting
  similar rhythmic behaviors.
\newblock {\em Current Biology}, 21:1036, 2011.

\bibitem{NSLGK012}
JM~Newcomb, A~Sakurai, JL~Lillvis, CA~Gunaratne, and PS~Katz.
\newblock Homology and homoplasy of swimming behaviors and neural circuits in
  the nudipleura (mollusca, gastropoda, 
  opistho-branchia).
\newblock {\em Proc. Natl. Acad. Sci.}, 109(1):10669--76, 2012.

\bibitem{sck2014}
A~Sakurai, CA~Gunaratne, and PS~Katz.
\newblock Two interconnected kernels of reciprocally inhibitory interneurons
  underlie alternating left-right swim motor pattern generation in the mollusk
  Melibe leonina.
\newblock {\em J. {N}europhysiology}, 112(6):317--1328, 2014.

\bibitem{sakurai2015phylogenetic}
A~Sakurai and PS~Katz.
\newblock Phylogenetic and individual variation in gastropod central pattern
  generators.
\newblock {\em J. Comparative Physiology A}, 201(9):829--839, 2015.


\bibitem{sakurai2019command}
A~Sakurai and PS~Katz.
\newblock Command or obey? homologous neurons differ in hierarchical position
  for the generation of homologous behaviors.
\newblock {\em J. Neuroscience}, 39(33):6460--6471, 2019.

\bibitem{sakurai2022bursting}
A~Sakurai and PS~Katz.
\newblock Bursting emerges from the complementary roles of neurons in a
  four-cell network.
\newblock {\em J. {N}europhysiology}, 127(4):1054--1066, 2022.



\bibitem{SakuraiKatz2017}
A~Sakurai and PS~Katz.
\newblock Artificial Synaptic Rewiring Demonstrates that Distinct Neural Circuit Configurations
Underlie Homologous Behaviors \newblock {\em Current Biology}, 27(4):1721--1734, 2017.


\bibitem{SakuraiKatz2016}
A~Sakurai and PS~Katz.
\newblock  The central pattern generator underlying swimming in \emph{Dendronotus iris}: A simple half-center network oscillator with a twist.
 \newblock {\em Journal of Neurophysiology}, 116(4):1728--1742, 2016.


\bibitem{Newcomb2012}
JM~Newcomb, A~Sakurai, J~LLillvis, CA~Gunaratne, and PS~Katz.
\newblock Homology and homoplasy of swimming behaviors and neural circuits in the Nudipleura (Mollusca, Gastropoda, Opisthobranchia).
\newblock {\em Proceedings of the National Academy of Sciences USA}, 109(Suppl 1):10669--10676, 2012.

\bibitem{zucker2002short}
RS~Zucker and WG~Regehr.
\newblock Short-term synaptic plasticity.
\newblock {\em Annual Review of Physiology}, 64:355--405, 2002.

\bibitem{scully2024pairing}
J~Scully, J~Bourahmah, D~Bloom, and AL~Shilnikov.
\newblock Pairing cellular and synaptic dynamics yields flexible network oscillators.
\newblock {\em Frontiers in Network Physiology}, 4:1397151, 2024.

\bibitem{jackman2017mechanisms}
SL~Jackman and WG~Regehr.
\newblock The mechanisms and functions of synaptic facilitation.
\newblock {\em Neuron}, 94(3):447--464, 2017.

\bibitem{klyachko2006molecular}
VA~Klyachko and CF~Stevens.
\newblock Molecular mechanisms of short-term plasticity: Role of synapsin phosphorylation in augmentation and potentiation of spontaneous glutamate release.
\newblock {\em Frontiers in Synaptic Neuroscience}, 8:12, 2006.

\bibitem{NeherBrose2018}
E~Neher and N~Brose.
\newblock Dynamically primed synaptic vesicle states: Key to understand synaptic short-term plasticity.
\newblock {\em Neuron}, 100(6):1283--1291, 2018.

\bibitem{plant81}
RE~Plant and M~Kim.
\newblock Bifurcation and resonance in a model for bursting nerve cells.
\newblock {\em Journal of Mathematical Biology}, 11:15--32, 1981.

\bibitem{destexhe1994synthesis}
Destexhe, A., Mainen, Z.F. and Sejnowski, T.J.
\newblock Synthesis of models for excitable membranes, synaptic transmission and neuromodulation using a common kinetic formalism.
\newblock {\em Journal of computational neuroscience}, 1(3):195--230, 1994.

\end{thebibliography}

\end{document} 











\bibitem{CPG}
A~Selverston, editor.
\newblock {\em Model Neural Networks and Behavior}.
\newblock Springer, Berlin, 1985.


\bibitem{KatzHooper07}
PS~Katz and SL~Hooper.
\newblock Invertebrate central pattern generators.
\newblock In G~North and RR~Greenspan, editors, {\em Invertebrate
  Neurobiology}. Cold Spring Harbor Laboratory Press, NY, New York, 2007.
  
  \bibitem{CALABRESE201668}
 RL Calabrese, and BJ~Norris and A.~Wenning.
\newblock Microcircuit computation and evolution.
\newblock {\em Current Opinion in Neurobiology},
  41:68-77, 2016. 
  
 \bibitem{Katz2016}
PS~Katz.
\newblock Evolution of central pattern generators and rhythmic behaviours.
 {\em Philos Trans R Soc Lond B Biol Sci}, 371:20150057, 2016. 
  

  \bibitem{MKM2022}
 E~Marder, S. Kedia and EO~Morozova.
\newblock New insights from small rhythmic circuits.
\newblock {\em Current Opinion in Neurobiology},
  76:102610, 2022.



%\bibitem{M12}
%E~Marder.
%\newblock Neuromodulation of neuronal circuits: back to the future.
%\newblock {\em Neuron}, 76:1, 2012.

\bibitem{NW02}
WH~Watson and JM~Newcomb.
\newblock Modulation of swimming in the gastropod Melibe leonina by nitric
  oxide.
\newblock {\em J. Exp. Biology}, 205(3):397, 2002.

\bibitem{WNT02}
WH~Watson, JM~Newcomb, and S~Thompson.
\newblock Neural correlates of swimming behavior in Melibe leonina.
\newblock {\em Biol. Bulletin}, 203(2):152, 2002.

\bibitem{TW05}
S~Thompson and WH~Watson.
\newblock Central pattern generator for swimming in Melibe.
\newblock {\em J. Exp. Biology}, 208(7):1347, 2005.


\bibitem{selverston1976stomatogastric}
AI~Selverston, DF~Russell, JP~Miller, and DG~King.
\newblock The stomatogastric nervous system: structure and function of a small
  neural network.
\newblock {\em Progress in neurobiology}, 7:215--289, 1976.

\bibitem{stein1999neurons}
PSG Stein, DG~Stuart, S~Grillner, and AI~Selverston.
\newblock {\em Neurons, networks, and motor behavior}.
\newblock MIT press, 1999.

\bibitem{selverston2013model}
AI~Selverston.
\newblock {\em Model neural networks and behavior}.
\newblock Springer Science \& Business Media, 2013.

\bibitem{Marder1994752}
E~Marder.
\newblock Invertebrate neurobiology: Polymorphic neural networks.
\newblock {\em Current Biology}, 4(8):752 -- 754, 1994.

%\bibitem{prinz2003functional}
%AA~Prinz, V~Thirumalai, and E~Marder.
%\newblock The functional consequences of changes in the strength and duration
%  of synaptic inputs to oscillatory neurons.
%\newblock {\em J. Neuroscience}, 23(3):943--954, 2003.

\bibitem{prinz2004similar}
AA~Prinz, D~Bucher, and E~Marder.
\newblock Similar network activity from disparate circuit parameters.
\newblock {\em Nature Neuroscience}, 7(12):1345--1352, 2004.

\bibitem{M09}
JM~Goaillard, AL~Taylor, D~Schulz, and E~Marder.
\newblock Multi-dimensional measurements to understand the functional
  consequences of animal-to-animal variation in circuit parameters.
\newblock {\em Nature Neuroscience}, 12:1424--1430, 2009.

\bibitem{Alacam2015}
D~Ala\c{c}am and AL~Shilnikov.
\newblock Making a swim central pattern generator out of latent parabolic
  bursters.
\newblock {\em International Journal of Bifurcation and Chaos}, 25(07):1540003,
  2015.



\bibitem{Perkel-Mulloney-74}
DH~Perkel and B~Mulloney.
\newblock Mechanism of post-inhibitory rebound in molluscan neurons.
\newblock {\em Science}, 185(4146):181--183, 1974.

\bibitem{SKM94}
F~Skinner, N~Kopell, and E~Marder.
\newblock Mechanisms for oscillation and frequency control in networks of
  mutually inhibitory relaxation oscillators.
\newblock {\em J. Comput. Neuroscience.}, 1:69, 1994.

\bibitem{SZVC99}
FK~Skinner, L~Zhang, JL~Perez~Velazquez, and PL~Carlen.
\newblock Bursting in inhibitory interneuronal networks: A role for
  gap-junctional coupling.
\newblock {\em J. Neurophysiol}, 81:1274, 1999.

\bibitem{Kopell-Ermentrout-02}
N~Kopell and G~Ermentrout.
\newblock {Mechanisms of phase-locking and frequency control in pairs of
  coupled neural oscillators}.
\newblock {\em Handbook on {D}ynamical {S}ystems}, 2:3--54, 2002.

\bibitem{Angstadt-Grassmann-Theriault-Levasseur-05}
JD~Angstadt, JL~Grassmann, KM~Theriault, and SM~Levasseur.
\newblock {Mechanisms of post-inhibitory rebound and its modulation by
  serotonin in excitatory swim motor neurons of the medicinal leech}.
\newblock {\em Journal of Comparative Physiology A-Neuroethology, Sensory,
  Neural and Behavioral Physiology}, 191(8):715--732, 2005.
  
  \bibitem{Wang-Rinzel-92}
XJ~Wang and J~Rinzel.
\newblock Alternating and synchronous rhythms in reciprocally inhibitory model
  neurons.
\newblock {\em {J. Neural Computations}}, {4}:84--97, {1992}.

\bibitem{Matveev-Bose-Nadim-Farzan-07}
V~Matveev, A~Bose, and F~Nadim.
\newblock {Capturing the bursting dynamics of a two-cell inhibitory network
  using a one-dimensional map}.
\newblock {\em {J. Computational Neuroscience}}, {23}({2}):{169--187},
  {OCT} {2007}.

\bibitem{Silverston2009}
A~Szucs, R~Huera, MI~Rabinovich, and AL~Silverston.
\newblock Robust microcircuit synchronization by inhibitory connections.
\newblock {\em Neuron}, 78(61):439--453, 2009.

\bibitem{Daun-Rubin-Rybak-09}
S~Daun, JE~Rubin, and IA~Rybak.
\newblock Control of oscillation periods and phase durations in half-center
  central pattern generators: a comparative mechanistic analysis.
\newblock {\em J. Computational Neuroscience}, 27(1):3--36, 2009.

\bibitem{nagornov2016mixed}
R~Nagornov, G~Osipov, M~Komarov, A~Pikovsky, and AL~Shilnikov.
\newblock Mixed-mode synchronization between two inhibitory neurons with
  post-inhibitory rebound.
\newblock {\em Communications in Nonlinear Science and Numerical Simulation},
  36:175--191, 2016.
  
\bibitem{SSN1996}
AA~Sharp, FK~Skinner, and E~Marder.
\newblock Mechanisms of oscillations in dynamics clap constructed two cell half-center circuits.
\newblock {\em J Neurophysiology},
  76(2):867--883, 1996


\bibitem{Jalil-Belykh-Shilnikov-10}
S~Jalil, IV~Belykh, and AL~Shilnikov.
\newblock {Fast reciprocal inhibition can synchronize bursting neurons}.
\newblock {\em {Physical Review E}}, {81}({4, 2}), {APR} {2010}.


\bibitem{Shilnikov2008b}
AL~Shilnikov, R~Gordon, and IV~Belykh.
\newblock Polyrhythmic synchronization in bursting networking motifs.
\newblock {\em Chaos}, 18(3):037120, Sep 2008.

\bibitem{brown}
TG~Brown.
\newblock The intrinsic factors in the act of progression in the mammal.
\newblock {\em London B Biol. Society}, 84(572):308--319, 1911.

\bibitem{wilson-cowen}
HR~Wilson , JD~Cowan.
\newblock{Excitatory and inhibitory interactions in localized populations of model neurons}
\newblock{\em{Biophysical journal}} 12(1):1-24, Jan 1972.


\bibitem{pinging}
P~Tiesinga, TJ~Sejnowski
\newblock{Cortical enlightenment: are attentional gamma oscillations driven by ING or PING?}
\newblock{\em{Neuron}}, 63(6):727-32, Sep 2009.

\bibitem{ping-b-wang}
G~Buzsáki, XJ~Wang.
\newblock{Mechanisms of gamma oscillations.}
\newblock{\em{ Annual review of neuroscience}}, 35:203-25. Jul 2012.

\bibitem{ping-kopell}
C~Burgers, S~Epstein, NJ~Kopell.
\newblock{Background gamma rhythmicity and attention in cortical local circuits: a computational study}
\newblock{\em{Proceedings of the National Academy of Sciences.}}, 102(19):7002-7, May 2005.

\bibitem{spindle-rinzel}
D~Golomb, XJ~Wang, J~Rinzel.
\newblock{Synchronization properties of spindle oscillations in a thalamic reticular nucleus model.}
\newblock{\em{Journal of neurophysiology}}, 72(3):1109-26, Sep 1994.
\bibitem{working-memory-wang}
N~Brunel, XJ~Wang.
\newblock{Effects of neuromodulation in a cortical network model of object working memory dominated by recurrent inhibition.}
\newblock{\em{Journal of computational neuroscience}}, 11:63-85, Jul 2001

%\bibitem{static-rubin-terman}
%J~Rubin, D~Terman.
%\newblock{Analysis of clustered firing patterns in synaptically coupled networks of oscillators.}
%\newblock{\em{Journal of mathematical biology}}, 41:513-45, Dec 2000.
%
%\bibitem{parkinsons-rubin}
%J~Rubin, D~Terman.
%\newblock{High frequency stimulation of the sub-thalamic nucleus eliminates pathological thalamic rhythmicity in a computational model.}
%\newblock{\em{Journal of computational neuroscience.}}, 16:211-35, May 2004.


\bibitem{rubin-06}
J~Rubin.
\newblock{Bursting induced by excitatory synaptic coupling in nonidentical conditional relaxation oscillators or square-wave bursters.}
\newblock {\em{Physical Review E}}, 74(2):021917, 2006.

\bibitem{plant75}
RE~Plant and M~Kim.
\newblock On the mechanism underlying bursting in the Aplysia abdominal
  ganglion R15 cell.
\newblock {\em Math. Biosciences}, 26:357--375, 1975.

\bibitem{plant76}
RE~Plant and M~Kim.
\newblock Mathematical description of a bursting pacemaker neuron by a
  modification of the Hodgkin-Huxley equations.
\newblock {\em Biophysics Journal}, 16:227--244, 1976.

\bibitem{plant81}
RE~Plant.
\newblock Bifurcation and resonance in a model for bursting nerve cells.
\newblock {\em J. Math. Biol}, 11:15--32, 1981.

\bibitem{RP1985}
AR~Rittenhouse and CH~Price.
\newblock Peripheral axons of the parabolic burster neuron R15.
\newblock {\em Brain Research}, 333:330--335, 1985.

\bibitem{LL88}
ES~Levitan and IB~Levitan.
\newblock Serotonin acting via cyclic amp enhances both the hyper-polarizing and
  depolarizing phases of bursting pacemaker activity in the {\em Aplysia} neuron
  {R}15.
\newblock {\em J Neuroscience}, 8:1152--1161, 1988.

\bibitem{Rinzel1985}
J~Rinzel.
\newblock Bursting oscillations in an excitable membrane model.
\newblock {\em Lecture Notes in Mathematics}, 1151:304--316, 1985.


\bibitem{Rinzel87b}
J~Rinzel.
\newblock A formal classification of bursting mechanisms in excitable systems. \newblock in {\em Proceedings of the International Congress of Mathematicians}, AMS, 1578--1593, 
\newblock 1987.

\bibitem{RL86}
J~Rinzel and YS~Lee.
\newblock {\em On different mechanisms for membrane potential bursting},
  volume~66.
\newblock In: Othmer, H.G. (ed.) Nonlinear Oscillations in Biology and
  Chemistry: Lecture Notes in Biomathematics, Springer, Berlin, 1986.

\bibitem{RL87}
J~Rinzel and YS~Lee.
\newblock Dissection of a model for neuronal parabolic bursting.
\newblock {\em J. Mathematical Biology}, 25(6):653--675, 1987.


\bibitem{CCB91}
CC~Canavier, JW~Clark, and JH~Byrne.
\newblock Simulation of the bursting activity of neuron {R}15 in aplysia: Role
  of ionic currents, calcium balance, and modulatory transmitters.
\newblock {\em J. Neurophysiology}, 66(6):2107--2124., 1991.

\bibitem{bertram1993}
R~Bertran.
\newblock A computational study of the effects of serotonin on a molluscan
  burster neuron.
\newblock {\em Biological Cybernetics}, 69:257--267, 1993.

\bibitem{Canavier1993}
CC~Canavier, DA~Baxter, JW~Clark, and JH~Byrne.
\newblock Nonlinear dynamics in a model neuron provide a novel mechanism for
  transient synaptic inputs to produce long-term alterations of postsynaptic
  activity.
\newblock {\em J Neurophysiology}, 69(6):2252--2257, Jun 1993.

\bibitem{BCCBB95}
RJ~Butera, J~W~Jr. Clark, CC~Canavier, DA~Baxter, and JH~Byrne.
\newblock Analysis of the effects of modulatory agents on a modeled bursting
  neuron: Dynamic interactions between voltage and calcium dependent systems.
\newblock {\em J. Computational Neuroscience}, 2(1):19--44, 1995.

\bibitem{butera98}
R~Butera.
\newblock Multirhythmic bursting.
\newblock {\em Chaos}, 8(1):274--284, 1998.

\bibitem{R15}
FH~Sieling and R~Butera.
\newblock Aplysia {R}15 neuron.
\newblock {\em Scholarpedia}, 6(10):4181, 2011.





\end{document}